\section{Related Work}
\label{sec:related}

This section is written from the reference-to-claim matrix in
\texttt{paper/RELATED\_WORK\_MATRIX.md}, and every reference was verified
against publisher-deposited metadata (\texttt{paper/REFERENCE\_AUDIT.csv}).
Where a body of work contains exceptions to a generalisation we might
otherwise draw, the exception is stated and the generalisation weakened
accordingly.

\subsection{Event-triggered multi-agent communication}
\label{sec:rel_etc}

Event-triggered control replaces periodic sampling with a condition on
locally measurable error. The canonical formulation transmits when the
state has drifted past a threshold from the last transmitted value, with a
minimum inter-event time to exclude Zeno behaviour~\cite{tabuada2007event},
and the event-versus-periodic comparison goes back
further~\cite{astrom2002riemann}; \cite{heemels2012introduction} gives the
standard taxonomy separating event-triggered from self-triggered designs.
The multi-agent extension is
mature: \cite{dimarogonas2012distributed} establishes distributed
event-triggered control from local state error,
\cite{seyboth2013eventbased} gives event-based broadcasting for average
consensus with an explicit broadcast count, and \cite{nowzari2019eventtriggered}
surveys the area. Thresholds need not be static: dynamic triggering
introduces an auxiliary variable that modulates the threshold over
time~\cite{girard2015dynamic,yi2017distributed}, and recent surveys
document how far that idea has been
carried~\cite{chen2020howoften,ge2021dynamic,zhang2025overview}.

This machinery has already been applied to multi-UAV formations, and we
state that plainly because it bounds what we can claim. Event-based
formation control has been developed for directed
topologies~\cite{yin2023eventbased}, under communication delay and
external disturbances~\cite{ji2023dynamic}, for formation tracking with
unknown disturbances~\cite{chen2024distributed}, and for swarm target
fencing with connectivity and collision
constraints~\cite{yang2025fencing}. \textbf{The gap this paper addresses is
therefore not the absence of event-triggered UAV formation control.}

Most representative designs above trigger on a local error --- the
difference between an agent's current state and the state it last sent ---
and are blind to what the receiver actually holds. There are important
exceptions to that generalisation. Wang et al. combine AoI with an
event-triggered distributed Kalman consensus filter over a wireless sensor
network~\cite{wang2021freshness}, and a recent consensus design triggers on
either state error or a sender-held AoI bound~\cite{lin2026cooperative}.
Neither design infers per-receiver freshness from delayed acknowledgements:
the latter explicitly resets its age variable when the sender broadcasts.
Thus ``AoI-aware event triggering'' is not itself the gap addressed here.

\subsection{Age of information and freshness in networked control}
\label{sec:rel_aoi}

Age of information supplies the missing quantity. Introduced as a freshness
metric for status-update systems~\cite{kaul2012realtime}, it measures how
old a receiver's copy is rather than how much traffic crossed the link, and
optimising it is not the same as transmitting as often as
possible~\cite{sun2017update}. The area is now
surveyed at length~\cite{yates2021aoisurvey,yang2025aoiperspective}, with
extensions to multi-node dissemination~\cite{kaswan2025gossip}.

Two results from this literature shaped our design directly. One is that age
alone is a poor proxy for control cost: \cite{ayan2019aoivoi} shows that
scheduling on value-of-information rather than age changes the outcome for
networked control loops. Our policy accordingly never lets age trigger a
transmission by itself --- age only lowers the bar that a genuine state
change must clear. Second, going beyond pure age is not new:
\cite{rajaraman2021notjustage} already combines age with a
content-quality term, so combining innovation-like content with freshness
already has precedent.

Closest to our combination is \cite{mamduhi2020freshness}, which co-designs
control and networking policy for multi-loop systems using a mixture of
age-of-information and event-based metrics. The setting, however, is a
network scheduler arbitrating among independent control loops. There is no
sender that maintains its own estimate of a particular receiver's
freshness, and no distinction between traffic that carries new information
and traffic that repeats it.

\subsection{Feedback- and ACK-aware communication and scheduling}
\label{sec:rel_ack}

A sender cannot observe the age of information at a receiver; it can only
infer it from what comes back. That inference is established in the
AoI-scheduling literature: \cite{ceran2019average} drives AoI-optimal
transmission from acknowledgement/negative-acknowledgement feedback with
hybrid automatic repeat request (ARQ) under a resource constraint, with no
oracle knowledge of future outcomes. More directly, Tahir et al. give
decentralised sensor agents delayed acknowledgements over a duplex channel
and use a particle filter to maintain a belief over receiver AoI in a
multi-agent transmission policy~\cite{tahir2024collaborative}. We therefore
do not claim either ACK-based causal freshness inference or its multi-agent
use as novel in itself.

Tahir et al. is the closest feedback-side near-neighbour found in our
independent re-check. Its objective and action depend on AoI and channel
load; it has no physical state-innovation term, no UAV formation-control
loop, and no semantic split between new information and refresh. It also
explicitly permits more than one message from an agent to be in flight.
Those differences, rather than the mere presence of ACKs, define our
narrower mechanism claim.

The contrast class is centralised and index-based: \cite{tang2022whittle}
schedules to minimise an age cost with a Whittle-index policy. Such
policies evaluate globally and periodically, where our policy decides
locally and asynchronously. Radio resources being an explicit design
variable in wireless control~\cite{gatsis2014optimal,park2018wireless} is
why we price the reverse channel rather than treating acknowledgements as
free, and delay and dropout affecting closed-loop
behaviour~\cite{walsh2002stability,hespanha2007survey} is why our network
model carries loss, delay, jitter and out-of-order delivery together.

Event-triggered mechanisms continue to be surveyed
actively~\cite{li2026impulsive}, which is one reason we position against
the current state of the field rather than against its founding papers
alone.

One difference in mechanism is worth noting here. \cite{ceran2019average}
recovers from loss with an explicit ARQ scheme. Our policy introduces no
retransmission timeout, no round-trip-time estimator and no window: an
unacknowledged packet remains outstanding, outstanding packets forbid
refresh traffic, and a max-silence backstop provides the floor.

\subsection{UAV-swarm communication and control co-design}
\label{sec:rel_uav}

The constraints that make fewer transmissions worth paying accuracy for
are documented
in the UAV networking literature: limited bandwidth, contested airtime and
reliability limits~\cite{gupta2016survey,zeng2016wireless},
with tutorial treatments of the design
space~\cite{mozaffari2019tutorial,zeng2019accessing} and evidence that
swarm architectures degrade as neighbour count
grows~\cite{campion2019uavswarm}. Age of information is used actively as a
design objective in this setting~\cite{amodu2023aoiuav}.

It has also reached hardware. \cite{tripathi2023wiswarm} demonstrates
AoI-based wireless networking for collaborative teams of UAVs on a real
platform, using a centralised scheduler to allocate the channel. That work
is both the nearest neighbour to ours on the application side and the
honest benchmark for our principal limitation: \textbf{our results are
simulation only}. Deployed aerial swarms
exist~\cite{vasarhelyi2018optimized,zhou2022swarm}, and a simulation-only
contribution must be positioned against them rather than beside them.

That communication constraints fundamentally limit achievable control
performance is long
established~\cite{tatikonda2004control,nair2007feedback}, and it is the
reason sending fewer packets is worth paying accuracy for at all.

On the control side we deliberately change nothing. The formation law is a
standard consensus-style interaction over a neighbour
graph~\cite{olfatisaber2006flocking,oh2015survey}, and the 6-DOF follower
model uses the usual cascaded attitude-thrust structure and the
acceleration-to-attitude reference
mapping~\cite{lee2010geometric,mellinger2011minimum}. Freezing the
controller is what allows the communication policy to be the only variable
--- and it is also why several failures we report are attributed to the
controller rather than to our contribution. Joint controller-and-trigger
design~\cite{molin2013optimality} and state-based
scheduling~\cite{ramesh2013design} are established alternatives that we do
not pursue; the communication-versus-control-cost trade-off they address is
studied explicitly in~\cite{demirel2017tradeoff}, which is our reason for
reporting cost under several models rather than one.

\subsection{Gap addressed by this work}
\label{sec:rel_gap}

Reading across the four subsections:

\begin{itemize}
\item event-triggered multi-agent and UAV-formation designs establish local
      state-error triggering, while AoI-aware consensus filters establish
      freshness-aware triggering; neither line supplies the full mechanism
      studied here (\S\ref{sec:rel_etc});
\item AoI-aware control designs model staleness, but usually place the
      decision in a scheduler or reset age at transmission rather than
      infer per-receiver freshness from delayed feedback
      (\S\ref{sec:rel_aoi});
\item ACK-driven causal receiver-AoI inference already exists in both
      single-link scheduling and a decentralised multi-agent sensor policy,
      without a physical state-innovation trigger or the new-versus-refresh
      semantic split (\S\ref{sec:rel_ack});
\item UAV-swarm work supplies the constraints, and has reached hardware
      with a \emph{centralised} AoI scheduler (\S\ref{sec:rel_uav}).
\end{itemize}

The combination that remains is a \emph{distributed sender in a formation}
that estimates its own receiver's freshness \emph{only} from cumulative
acknowledgements, uses that estimate to modulate a \emph{state-innovation}
threshold, and separates \emph{new information from refresh} so that a
repetition cooldown governs only repetition. That combination is what this
paper addresses.

We are explicit about the strength of this statement. Each individual
ingredient is precedented. \textbf{We did not identify prior work combining}
these elements in this way --- a statement about our documented search, not
a proof of absence. The search boundary, the nearest neighbours with a per-element
comparison, and the conditions that would overturn the assessment are
recorded in \texttt{paper/NOVELTY\_GAP\_REVIEW.md}. No priority claim is
made.
